\begin{description}
\item[(T08.1)] Circular trajectories have $r = $ constant. Hence, the radial
  velocity $v_{\mathrm{r}} = 0$ \hfill(1.0)

  and rate of change of radial component of velocity
  $\dfrac{dv_{\mathrm{r}}}{dt} = 0$. \hfill(1.0)

  From the third equation,
  \[
    \frac{dv_{\mathrm{r}}}{dt} - \frac{L^2}{r^3} + \frac{3 G M L^2}{c^2 r^4}
    = 0,
  \]
  we get radius $r_{\mathrm{ph}}$ of circular orbits to be
  $r_{\mathrm{ph}} = 3GM/c^2$. \hfill(1.0)

  Then from the first equation, the impact parameter is found to be
  \[
    b_{\mathrm{ph}} = L/\sqrt{2E} = 3\sqrt{3}\,GM/c^2.
  \]
  \hfill(1.0)

\item[(T08.2)] Time taken to complete one circular orbit is
  \begin{align*}
    T_{\mathrm{ph}} &= 2\pi r_{\mathrm{ph}}/c = 6\pi G M/c^3
      \tag*{(1.0)}\\
      &= 6.0 \times 10^{5}\,\mathrm{s}
      \tag*{(1.0)}
  \end{align*}

\item[(T08.3a)] $r_\alpha$ corresponds to the radius where
  $V_{\mathrm{eff}} = \dfrac{L^2}{2r^2}\left(1 - \dfrac{2GM}{c^2 r}\right)$ is
  equal to zero. Hence $r_\alpha = 2GM/c^2 = R_{\mathrm{SC}}$. \hfill(1.0)

  Solving for $\dfrac{dV_{\mathrm{eff}}}{dr} = 0$ gives
  $r_\beta = 3GM/c^2 = r_{\mathrm{ph}}$. \hfill(1.0)

\item[(T08.3b)] The radius of the position of peak of $V_{\mathrm{eff}}(r)$ is
  the smallest possible turning point radius $r_{\mathrm{t}}$ for the light
  trajectory such that it can escape to infinity. \hfill(1.0)

  Hence, the least smallest possible turning point radius $r_{\mathrm{t}}$ is
  same as the radius $r_{\mathrm{ph}}$ of the circular orbits.
  $r_{\mathrm{t}} = 3GM/c^2$ \hfill(1.0)

  Since $r_{\mathrm{t}}$ is equal to $r_{\mathrm{ph}}$, the corresponding
  impact parameter is same as calculated in part one
  $b_{\mathrm{min}} = 3\sqrt{3}\,GM/c^2$. \hfill(1.0)

\item[(T08.4)] Since it is given that $r_{\mathrm{app}} \approx b$, the
  apparent radii $r_{\mathrm{inner}}$ and $r_{\mathrm{outer}}$ are related to
  the respective impact parameters $b_{\mathrm{inner}}$ and
  $b_{\mathrm{outer}}$ of the corresponding light rays.

  To figure out the $b_{\mathrm{inner}}$ and $b_{\mathrm{outer}}$, one needs to
  choose the correct corresponding $r_{\mathrm{actual}}$ in each case.

  The outer ring radius $r_{\mathrm{outer}}$ would correspond to the light ray
  which is emitted from the outer-edge of the accretion disk. Hence the
  $r_{\mathrm{outer}}$ is determined by the formula given
  $r_{\mathrm{outer}} \approx b_{\mathrm{outer}} \approx a_{\mathrm{outer}}
  \left(1 + R_{\mathrm{SC}}/a_{\mathrm{outer}}\right)^{1/2}$. \hfill(1.0)

  Substituting $a_{\mathrm{outer}} = 10 R_{\mathrm{SC}}$ we get
  $r_{\mathrm{outer}} = 1.3 \times 10^{3}$\,au. \hfill(1.0)

  For the inner ring radius $r_{\mathrm{inner}}$, the light ray corresponding
  to the smallest turning point $r_{\mathrm{t}}$ is the light ray which goes
  closest to the BH and returns back to infinity.

  Thus one should choose the impact parameter $b_{\mathrm{inner}}$ to be that
  of the circular orbit. Then,
  $r_{\mathrm{inner}} \approx b_{\mathrm{inner}} = b_{\mathrm{ph}}
  = 3\sqrt{3}\,GM/c^2$. \hfill(2.0)

  Hence $r_{\mathrm{inner}} = 3.3 \times 10^{2}$\,au. \hfill(1.0)

\item[(T08.5a)] The source is a point source.

  For the image, one has to look at all the possible trajectories light that
  can reach the observer starting from point Z in the same plane as the
  observer, the black hole and point Z.

  The shortest path is the direct path, without orbiting the black hole,
  connecting Z and the observer. This will be also be bent due to gravity. The
  next shortest path could be the one which goes half around the black hole in
  the anti-clockwise sense in the picture. Another one is the one with one and
  half revolution around the black hole in the anti-clockwise sense in the
  picture. And so on, paths with additional one more revolution around the
  black hole than the previous one.

  The above is also true for trajectories in clockwise sense of rotation
  around the black hole. Therefore, number of possible paths for light to
  travel from Z to the observer is (D) Greater than two. \hfill(2.0)

\item[(T08.5b)] One has to consider the time difference in the arrival times
  due to the different path lengths. From Part (T08.2), it takes around
  $6.0 \times 10^{5}$\,sec which is 166.7\,h for light to circularly orbit the
  supermassive black hole of mass $M = 6.5 \times 10^{9}\,\mathrm{M_\odot}$.
  Hence, the difference in the arrival times of light along these trajectories
  would be at least around half of 166.7\,h approximately, that is, around
  83.3\,h. Each of these signals would last for only the duration of the
  blast, i.e., 5\,s.

  Since the long exposure picture is only for 60\,s, either one or no image
  will be captured, depending on whether the exposure time encompasses the
  short window of the arrival of one of the signals and 5\,s thereafter.
  Therefore, the number of images in a picture would be (A) At most one.
  \hfill(2.0)
\end{description}
